% =========================
% Cas pratique 4
% =========================

\section[Cas 4]{Cas pratique 4 : Recevabilité en
contrefaçon}

%--------------------------
\begin{frame}{Question 1}

\begin{longtblr}{colspec = {X[3,l] X[7,l] X[5,l]}, rowsep = 1pt}

\textbf{personne}
 & \textbf{condition?}
 & \\
\hline
\ok{déposant} et/ou \mydots
 & -- pas de condition
 & L615-2 al. 1 \\

\ok{co-pro} et/ou \mydots
 & -- notification aux copropriétaires
 & L613-29 \\

\ok{cessionnaire} et/ou \mydots
 & -- cession inscrite au RNB pour cessionnaire par transmission à titre particulier ou par transmission à titre universel
(p. ex: fusion, absorption, héritage)
 & L613-9 al. 1 \\

 & -- cession inscrite au RNB en procédure: régularisation rétroactive jusqu'au jour de la cession
 & Pourvoi n°D22-22.999 \\
\end{longtblr}

\end{frame}

\begin{frame}{Question 1}

\begin{longtblr}{colspec = {X[3,l] X[7,l] X[5,l]}, rowsep = 1pt}
\ok{licencié exclusif} et/ou \mydots
 & -- (i) notification au breveté et (ii) pas d'interdiction d'engager une action en cf.
 pour action en cessation de l'atteinte et action en réparation
 & L615-2 al. 2 \\

\ok{licencié non exclusif} et/ou \mydots
 & -- (i) notification au breveté et (ii) autorisation expression d'engager une action en contrefaçon
 %pour action en réparation de son préjudice ou intervention dans l'instance engagée par le breveté
 %pour réparation de son préjudice propre
 & L615-2 al. 6 \\

\ok{licencié}
 & -- pas de condition d'inscription au RNB pour se joindre à l'action du breveté
pour réparation de son préjudice propre
 & L613-9 al. 3\\
% &
% - pour les marques/modèles: licencié non inscrit au RNB engage une action &
% Cour de Justice 22/06/2016\\

\nok{cédant}
 & -- réparation du préjudice lorsque cession des droits de poursuite antérieurs
 & \\

\end{longtblr}

\end{frame}
%--------------------------

%--------------------------
\begin{frame}{Question 2}

\begin{block}{}
Art. \textbf{L615-2} al. 1 CPI: L'action en contrefaçon est exercée par le titulaire du brevet.
\end{block}

\begin{itemize}
    \item En l'espèce, brevet A déposé par la société France Tube.
    Pas de condition pour le déposant.
    \item En conclusion, France Tube est recevable à agir en contrefaçon du brevet A.
\end{itemize}

\begin{block}{}
Art. \textbf{L613-9} al. 1 CPI: Tous les actes transmettant ou modifiant 
les droits attachés à une demande de brevet ou à un brevet doivent, 
pour être opposables aux tiers, être inscrits sur un registre, 
dit Registre national des brevets, tenu par l'Institut national de la propriété industrielle.
\end{block}

\begin{itemize}
    \item En l'espèce, brevet B déposé par M. Gendre dirigeant de France Tube.
    \item En conclusion, France Tube est recevable à agir en contrefaçon du brevet A à
    condition d'inscrire la cession au RNB.
\end{itemize}

\end{frame}
%--------------------------

%--------------------------
\begin{frame}{Question 3}

\begin{block}{}
Art. \textbf{L615-2} al. 2 CPI: Sauf stipulation contraire du contrat de licence,
elle est également ouverte au titulaire d'une licence exclusive à condition,
à peine d'irrecevabilité, d'informer au préalable le titulaire du brevet.

\end{block}

\begin{itemize}
    \item En l'espèce, France Tube est licenciée exclusive du brevet C.
    \item En conclusion, France Tube est recevable à agir en contrefaçon du brevet C à
    condition de notifier le breveté et de ne pas avoir d'interdiction d'engager
    une action en contrefaçon.
\end{itemize}

\end{frame}
%--------------------------

%--------------------------
\begin{frame}{Question 4}

\begin{block}{}
Art. \textbf{47}(1) AJUB dispose que le titulaire du brevet européen a qualité pour agir
    en contrefaçon devant la Juridiction Unifiée du Brevet. \\
Art. \textbf{47}(2) AJUB Le licencié exclusif peut également agir, saus si l'accord
de licence en dispose autrement et à condition d'informer le titulaire du brevet.
\end{block}

\begin{itemize}
    \item En l'espèce, France Tube aurait souhaité agir devant la JUB car ses brevets EP
    n'ont pas fait l'objet d'une déclaration de dérogation (opt-out).
    \item En conclusion, devant la JUB, France Tube est recevable à agir en contrefaçon du brevet A et B.
    France Tube est recevable à agir en contrefaçon du brevet C à
    condition de notifier le breveté et de ne pas avoir d'interdiction d'engager
    une action en contrefaçon.
\end{itemize}

\begin{itemize}
    \item Les règles JUB sont proches de celles du droit français.
\end{itemize}

\end{frame}
%--------------------------